\section*{अध्याय \ref{ch:b1:alcohol-activation} --- ऐल्कोहॉल: OH समूह का सक्रियण}

\begin{solution}{exo:b1:alcohol-activation:1}
\ce{2CH3CH2OH + 2Na -> 2CH3CH2O- + 2Na+ + H2};
\ce{CH3CH2OH + NaH -> CH3CH2O- + Na+ + H2};
\ce{CH3CH2OH + OH- <=> CH3CH2O- + H2O}, $K = 10^{14.0 - 15.9} = 10^{-1.9}$:
आंशिक।
\end{solution}

\begin{solution}{exo:b1:alcohol-activation:2}
1-मेथॉक्सीप्रोपेन; मेथॉक्सीएथेन; एथॉक्सीबेंज़ीन (फ़ीनॉक्साइड, जो \omterm{def:b1:alcohol-activation:alkoxide}{ऐल्कॉक्साइड} से
\omterm{def:b1:predominant-reaction:strong-acid}{दुर्बल क्षारक} है, फिर भी एक अच्छा \omterm{def:b1:electronic-effects:nucleophile}{नाभिकरागी} है)।
\end{solution}

\begin{solution}{exo:b1:alcohol-activation:3}
पिरिडीन में टोसिलन से \ce{CH3CH2CH2OTs} मिलता है; फिर
\[ \ce{CH3CH2CH2OTs + I- -> CH3CH2CH2I + TsO-}. \] प्रत्यक्ष रूप से, आयोडाइड को
\ce{OH-} को बाहर निकालना पड़ता: कोई अभिक्रिया नहीं। \omterm{def:b1:alcohol-activation:sulfonate}{टोसिलेट} में एक उत्कृष्ट \omterm{def:b1:electronic-effects:nucleophile}{निर्गामी समूह} है,
और परिस्थितियाँ न अम्लीय हैं न प्रबल क्षारकीय।
\end{solution}

\begin{solution}{exo:b1:alcohol-activation:4}
ब्यूटेन-2-ऑल: ब्यूट-1-ईन, \cip{Z}- और \cip{E}-ब्यूट-2-ईन; मुख्य
\cip{E}-ब्यूट-2-ईन। 2-मेथिलपेंटेन-2-ऑल: 2-मेथिलपेंट-2-ईन (मुख्य,
त्रिप्रतिस्थापित) और 2-मेथिलपेंट-1-ईन।
\end{solution}

\begin{solution}{exo:b1:alcohol-activation:5}
\ce{(CH3)2CHCH2-O-CH2CH3}: ब्रोमोएथेन के साथ सोडियम 2-मेथिलप्रोपेन-1-ओलेट,
या 1-ब्रोमो-2-मेथिलप्रोपेन के साथ सोडियम एथॉक्साइड। दोनों हैलाइड प्राथमिक हैं,
पर 1-ब्रोमो-2-मेथिलप्रोपेन अभिक्रिया करने वाले कार्बन के ठीक बगल में शाखित है,
जो SN2 को धीमा करता है और \omterm{def:b1:predominant-reaction:strong-acid}{प्रबल क्षारक} को विलोपन करने देता है (2-मेथिलप्रोपीन)।
पहला चयन बेहतर है।
\end{solution}

\begin{solution}{exo:b1:alcohol-activation:6}
\omterm{def:b1:alcohol-activation:sulfonate}{टोसिलेट} \cip{R} \omterm{def:b1:stereochemistry-in-depth:isomers}{विन्यास} बनाए रखता है; एथॉक्साइड के साथ SN2 उसे प्रतीपित
करता है: \cip{S}-2-एथॉक्सीऑक्टेन (OEt का क्रम पहला है, जैसे OTs का था)। सल्फ़्यूरिक अम्ल
और ऊष्मा के साथ: एक द्वितीयक \omterm{def:b1:electronic-effects:intermediates}{कार्बधनायन} के माध्यम से ऑक्टीनों तक विलोपन
(मुख्यतः \cip{E}-ऑक्ट-2-ईन), त्रिविम रसायन की हानि के साथ।
\end{solution}

\begin{solution}{exo:b1:alcohol-activation:7}
OH का प्रोटॉनन; तृतीयक \omterm{def:b1:electronic-effects:intermediates}{कार्बधनायन} तक जल की हानि (मंद);
\ce{Cl-} द्वारा ग्रहण: 2-क्लोरो-2-मेथिलप्रोपेन। तीव्र, क्योंकि तृतीयक
\omterm{def:b1:electronic-effects:intermediates}{कार्बधनायन} स्थायी है और क्लोराइड, जो अम्ल में अविलेय है, अलग हो जाता है।
\end{solution}

\begin{solution}{exo:b1:alcohol-activation:8}
\ce{CH3CH2OH + H+ -> CH3CH2OH2+}; एथेनॉल का दूसरा अणु कार्बन पर
पीछे से आक्रमण करता है और जल निकलता है (SN2), जिससे \ce{(CH3CH2)2OH+} बनता है, जो एक प्रोटॉन खोकर
एथॉक्सीएथेन देता है। प्रतिस्पर्धी अभिक्रिया: एक क्षारक (एथेनॉल, \ce{HSO4-}) \ce{CH3CH2OH2+} का एक $\beta$
प्रोटॉन लेता है जबकि जल निकलता है (E2): एथीन।
\end{solution}

\begin{solution}{exo:b1:alcohol-activation:9}
2-मेथिलप्रोपेन-2-ऑल (तृतीयक), जो अविलेय 2-क्लोरो-2-मेथिलप्रोपेन देता है।
\end{solution}

\begin{solution}{exo:b1:alcohol-activation:10}
प्रोटॉनन और जल की हानि एक चतुष्क कार्बन के बगल में एक द्वितीयक \omterm{def:b1:electronic-effects:intermediates}{कार्बधनायन}
देते हैं। उस कार्बन का एक मेथिल समूह अपने \omterm{def:b1:lewis-resonance-vsepr:lewis}{आबंधी युग्म} के साथ
धनायनी कार्बन पर खिसक जाता है: धन आवेश अब एक \omterm{def:b1:electronic-effects:carbon-class}{तृतीयक
कार्बन} पर है, जो अधिक स्थायी है। पड़ोसी कार्बन से एक प्रोटॉन की हानि
चतुष्प्रतिस्थापित 2,3-डाइमेथिलब्यूट-2-ईन देती है, जो सबसे स्थायी ऐल्कीन है।
\end{solution}

\begin{solution}{exo:b1:alcohol-activation:11}
\cip{R}-ब्यूटेन-2-ऑल से: \omterm{def:b1:alcohol-activation:sulfonate}{टोसिलेट} (\cip{R}), फिर सोडियम मेथॉक्साइड (SN2,
प्रतीपन): \cip{S}-2-मेथॉक्सीब्यूटेन। \cip{S}-ब्यूटेन-2-ऑल से:
\omterm{def:b1:alcohol-activation:alkoxide}{ऐल्कॉक्साइड} (NaH) बनाइए और आयोडोमेथेन मिलाइए: त्रिविमजनक C--O आबंध को
छुआ नहीं जाता, फिर से \cip{S}-2-मेथॉक्सीब्यूटेन।
\end{solution}

\begin{solution}{exo:b1:alcohol-activation:12}
$Q = [\text{ईथर}][\ce{H2O}]/[\ce{EtOH}]^2$। बनते ही जल हटाते रहने से
$[\ce{H2O}]$ और $Q$ छोटे रहते हैं: $Q < K$ और अभिक्रिया तब तक ईथर बनाती रहती है
जब तक ऐल्कोहॉल खप न जाए।
\end{solution}

\begin{solution}{pb:b1:alcohol-activation:1}
\textbf{1.} मेथिल हैलाइड के साथ सोडियम टर्ट-ब्यूटॉक्साइड; टर्ट-ब्यूटिल हैलाइड के साथ
सोडियम मेथॉक्साइड।
\textbf{2.} दूसरा: मेथॉक्साइड, जो \omterm{def:b1:predominant-reaction:strong-acid}{प्रबल क्षारक} है, तृतीयक हैलाइड से
विलोपन करता है, \ce{(CH3)3CBr + CH3O- -> (CH3)2C=CH2 + CH3OH + Br-}।
\textbf{3.} \ce{(CH3)3CO- + CH3I -> (CH3)3COCH3 + I-}: \omterm{def:b1:alcohol-activation:alkoxide}{ऐल्कॉक्साइड} ऑक्सीजन
मेथिल कार्बन पर पीछे से आक्रमण करता है जबकि आयोडाइड निकलता है (SN2)।
\textbf{4.} सोडियम (या सोडियम हाइड्राइड) के साथ: \ce{2(CH3)3COH + 2Na ->
2(CH3)3CO- + 2Na+ + H2}। हाइड्रॉक्साइड बहुत \omterm{def:b1:predominant-reaction:strong-acid}{दुर्बल क्षारक} है: $K = 10^{14.0 -
19.2} = 10^{-5.2}$।
\textbf{5.} जल \omterm{def:b1:alcohol-activation:alkoxide}{ऐल्कॉक्साइड} को प्रोटॉनित करेगा और हैलाइड पर आक्रमण करेगा।
\textbf{6.} नहीं: कोई \omterm{def:b1:stereochemistry-in-depth:stereocentre}{त्रिविमजनक केंद्र} नहीं।
\textbf{7.} \cip{R}-ब्यूटेन-2-ऑल + TsCl (पिरिडीन) $\to$ \cip{R}-ब्यूटेन-2-इल
\omterm{def:b1:alcohol-activation:sulfonate}{टोसिलेट}: धारण।
\textbf{8.} \cip{S}-2-मेथॉक्सीब्यूटेन, SN2 प्रतीपन द्वारा।
\textbf{9.} \cip{S}-ब्यूटेन-2-ऑल।
\textbf{10.} \cip{R}-2-मेथॉक्सीब्यूटेन: त्रिविमजनक कार्बन का C--O आबंध
बना रहता है; केवल O--C(मेथिल) बनता है।
\textbf{11.} E2 (मेथॉक्साइड \omterm{def:b1:predominant-reaction:strong-acid}{प्रबल क्षारक} है और कार्बन द्वितीयक है),
जो ब्यूटीन देता है; छोटे क्षारक, उत्कृष्ट \omterm{def:b1:electronic-effects:nucleophile}{निर्गामी समूह} और
मृदु ताप से सीमित।
\textbf{12.} प्रोटॉनित ऐल्कोहॉल आयनित होकर समतलीय द्वितीयक
\omterm{def:b1:electronic-effects:intermediates}{कार्बधनायन} बनाता है, जिसे मेथेनॉल दोनों फलकों पर ग्रहण करता है: रेसेमिक ईथर,
ब्यूटीनों के साथ।
\textbf{13.} OH का प्रोटॉनन, तृतीयक \omterm{def:b1:electronic-effects:intermediates}{कार्बधनायन} तक जल की हानि (मंद),
एक $\beta$ प्रोटॉन की हानि।
\textbf{14.} 2-मेथिलब्यूट-2-ईन (मुख्य, त्रिप्रतिस्थापित) और
2-मेथिलब्यूट-1-ईन।
\textbf{15.} तृतीयक \omterm{def:b1:electronic-effects:intermediates}{कार्बधनायन} बाधित है और ऐल्कोहॉल के दूसरे अणु के आक्रमण की
तुलना में कहीं अधिक आसानी से प्रोटॉन खो देता है।
\textbf{16.} एक उत्पाद हटाने के लिए: $Q$ छोटा रखने पर अभिक्रिया चलती रहती है,
और ऐल्कीन अम्ल के संपर्क में नहीं रहती।
\textbf{17.} \omterm{def:b1:electronic-effects:intermediates}{कार्बधनायन} तक एक ऊँचा पहला अवरोध (जल की हानि, दर-निर्धारक),
फिर ऐल्कीन तक एक छोटा अवरोध।
\textbf{18.} नहीं: प्राथमिक \omterm{def:b1:electronic-effects:intermediates}{कार्बधनायन} बहुत अस्थायी है; प्राथमिक ऐल्कोहॉल
प्रोटॉनित ऐल्कोहॉल पर E2 द्वारा, उच्चतर ताप पर, विलोपन करते हैं।
\textbf{19.} $10.0/74.0 = \qty{0.135}{mol}$।
\textbf{20.} $1.10 \times 0.135 \times 141.9 = \qty{21.1}{g}$।
\textbf{21.} $0.135 \times 88.0 = \qty{11.9}{g}$।
\textbf{22.} $0.75 \times 11.9 =$ \textbf{\qty{8.9}{g} MTBE}।
\end{solution}
